\scanpage{8}
\uline{Remarque} \uline{5 bis} \quad Avec les notations du lemme 5, et en présence des conditions (i) à (iv), la condition (v) équivaut à la condition suivante, qui se prête mieux à une vérification directe :

(v bis) La fonction $p\longmapsto\operatorname{rang}_{k(p)}F(k(p))$ \deleted{\unsure{sur|sera}} sur $S$ est constructible. \inserted{Il faut cependant supposer que $A$ satisfait à la condition : pour tout quotient intègre $B$ de $A$, $\unsure{\operatorname{Nor}|\operatorname{Nr}}(B)$ contient un ouvert non vide.}

Le fait que $(\mathrm v)\Longrightarrow(\mathrm v\ \mathrm{bis})$ est trivial.

En sens inverse, on utilise le fait que $F$ définit un ``syst. local'' de modules $\deleted{\unsure{M|Q}}P_p$. Dans le cas $A$ intègre, \deleted{l'hyp.} si le rang des $\deleted{M}P_p\otimes k(p)$ est constant et \unsure{égale|égal} à $n$ pour $p\in U$ (ouvert non vide), \deleted{et \unsure{x}} on en conclut que les $\deleted{M}P_p$ $(p\in U)$ sont alors \uline{libres} de rang $n$. D'ailleurs, \deleted{\unsure{de par}} comme $F$ commute aux $\varinjlim$, on voit qu'on peut supposer que les $\deleted{M}P_p$ sont tous $\subset N$, avec $N$ libre de rang $\unsure{n|N}$. Utilisant le fait que \deleted{\unsure{alors le}} $F(\deleted{\unsure{M}}N)$ est de type fini, on trouve aisément que pour tous les $p\in U$ de codim. 1, \uline{sauf un nombre fini}, on a $P_p=N_p$. Quitte à restreindre $U$, on peut supposer que ceci est vrai pour tous les $p$ de codim 1 de $U$. Mais alors, si on prend $U$ assez petit pour être normal, on conclut que $P_p=N_p$ pour \uline{tt} $p\in U$, \deleted{donc} d'où facilement que $N$ représente $F|U$.

% archive sheet number: 7
